\documentclass[a4paper]{article}
% Shared preamble for the English probability and statistics documents.
% Based on the original course preamble. Compile with pdfLaTeX.
\usepackage[top=1.5cm,bottom=2cm,left=2cm,right=2cm]{geometry}
\usepackage[T1]{fontenc}
\usepackage{lmodern}
\usepackage[utf8]{inputenc}
\usepackage[english]{babel}

\usepackage{amsmath}
\usepackage{amssymb}
\usepackage{amsthm}
\usepackage{mathrsfs}
\usepackage{mathtools}
\usepackage{etoolbox}
\usepackage{tikz}
\usetikzlibrary{shapes.multipart,arrows.meta}
\usepackage{xcolor}
\usepackage{enumerate}
\usepackage[shortlabels]{enumitem}
\usepackage{hyperref}
\usepackage{version}
\usepackage{bbold}
\usepackage{fontawesome5}
\usepackage{fancyhdr}
\lfoot{}
\usepackage{natbib}
\usepackage{graphicx}
\usepackage{pdfpages}
\usepackage{authblk}
    
\usepackage{blindtext}
\usepackage{multicol}
\usepackage{collcell}
\usepackage{wrapfig}
\usepackage{blkarray}
\usepackage{cancel}
\usepackage{appendix}
\usepackage{float}
\makeatletter
\input{tex/assets/slashbox.sty}
\makeatother
\hypersetup{unicode,colorlinks,citecolor=red,filecolor=black,
  linkcolor=red,urlcolor=cyan}
\setlength{\columnsep}{1cm}
\pagenumbering{gobble}

\def\numtext#1{%
  \ifcase#1\or one\or two\or three\or four\or five\or six\or seven\or eight\or nine\or ten\or
  eleven\or twelve\or thirteen\or fourteen\or fifteen\or sixteen\or seventeen\or eighteen\or nineteen\or twenty\or Lots\fi}

\newcommand{\Byte}[1]{%
  \def\nodecontents{}%
  \foreach[count=\l] \k in {#1} {\xappto{\nodecontents}{\noexpand\nodepart{\numtext{\l}}\k}}%
  \node[rectangle split,rectangle split parts=8,rectangle split horizontal,draw] at (2,2){\nodecontents};}

\newcommand{\D}{{\mathrm{d}}}
\newcommand{\cE}{{\mathcal{E}}}
\newcommand{\cG}{{\mathcal{G}}}
\newcommand{\cX}{{\mathcal{X}}}
\newcommand{\cY}{{\mathcal{Y}}}
\newcommand{\cB}{{\mathcal{B}}}
\newcommand{\cF}{{\mathcal{F}}}
\newcommand{\cO}{{\mathcal{O}}}
\newcommand{\cS}{{\mathcal{S}}}
\newcommand{\cT}{{\mathcal{T}}}
\newcommand{\given}{\, \big| \,}
\newcommand{\cL}{{\mathscr{L}}}
\newcommand{\cD}{{\mathcal{D}}}
\newcommand{\cA}{{\mathscr{A}}}
\newcommand{\cP}{{\mathscr{P}}}
\newcommand{\cR}{{\mathscr{R}}}
\newcommand{\A}{{\mathcal A}}
\newcommand{\one}{\mathbb{1}}
\newcommand{\mathbbm}[1]{\mathbb{#1}}

\providecommand{\C}{}
\renewcommand{\C}{\mathbb C}
\newcommand{\R}{\mathbb R}
\newcommand{\Q}{\mathbb Q}
\newcommand{\N}{\mathbb N}
\newcommand{\T}{\mathbb T}
\newcommand{\F}{\mathbb F}
\renewcommand{\P}{\mathbb P}
\newcommand{\Z}{\mathbb Z}
\newcommand{\E}{\mathbb {E}}
\newcommand{\Pbb}[1]{\P\left( #1\right)}
\newcommand{\Ebb}[1]{\E\left[ #1\right]}

\newcommand{\lb}{\left (}
\newcommand{\rb}{\right )}
\newcommand{\lbb}{\left [}
\newcommand{\rbb}{\right ]}
\newcommand{\labs}{\left |}
\newcommand{\rabs}{\right |}
\newcommand{\lbrb}[1]{\lb #1 \rb}
\newcommand{\lbbrbb}[1]{\lbb#1\rbb}
\newcommand{\labsrabs}[1]{\labs#1\rabs}
\newcommand{\lbbrb}[1]{\lbb#1\rb}
\newcommand{\lbrbb}[1]{\lb#1\rbb}
\newcommand{\langlerangle}[1]{\langle#1\rangle}
\newcommand{\lbcurly}{\{}
\newcommand{\rbcurly}{\}}
\newcommand{\curly}[1]{\lbcurly#1\rbcurly}
\newcommand{\abs}{\labsrabs}
\newcommand{\ind}[1]{1_{#1}}
\newcommand{\indep}{\perp \!\!\! \perp}
\renewcommand{\d}{\, \mathrm{d}}

\newtheorem{theorem}{Theorem}
\newtheorem{problemCnt}{Lemma}
\theoremstyle{definition}
\newtheorem{problem}[problemCnt]{Problem}
\newtheorem{problembook}[problemCnt]{\faBook \,  Problem}
\newtheorem{problempencil}[problemCnt]{\faPencil* \,  Problem}
\newtheorem{problemstar}[problemCnt]{\faStar\,  Problem}
\newtheorem{problemstarpencil}[problemCnt]{\faStar \, \faPencil* \,   Problem}
\newtheorem{probleminteresting}[problemCnt]{\faSmile[regular] \,   Problem}
\newtheorem{solution}{Solution}

\newcommand{\fwcell}[1]{\makebox[\arraystretch\normalbaselineskip]{$#1$}}
\newenvironment{enumeratei}{\begin{enumerate}[(i)]}{\end{enumerate}}
\newenvironment{enumeratea}{\begin{enumerate}[(a)]}{\end{enumerate}}
\newenvironment{enumeraten}{\begin{enumerate}[\upshape 1.]}{\end{enumerate}}

% Compatibility helpers needed by a few files in the combined archive.
\AtBeginDocument{\noindent{\footnotesize English translation prepared with ChatGPT (OpenAI, GPT-6), 26 October 2026, for translation of the Bulgarian text into English and LaTeX formatting.
The Bulgarian originals remain the reference; translation may contain errors.}\par\smallskip}
\newcommand{\MissingGraphic}[1]{%
  \begin{center}\fbox{\parbox{0.82\linewidth}{\centering Missing source image: \texttt{#1}}}\end{center}}



\setlength{\emergencystretch}{3em}
\pagenumbering{arabic}
\begin{document}
\begin{center}\Large Probability and Statistics, Computer Science\par\large Test 2, 9 June 2023: solutions\end{center}
\noindent\textit{English transcription of the supplied handwritten solution notes.}
\section*{Problem 1: leading digit 1}
For any distribution of $X$, let $A$ be the probability that the first nonzero decimal digit is 1. Sum the probabilities over $[10^k,2\cdot10^k)$, $k\in\mathbb Z$, including the corresponding negative intervals for negative values. Since accuracy 0.01 is requested, only a few terms may be needed. For Benford's law, the probability is $\log_{10}2\approx30.1\%$.
\subsection*{1. Standard normal}
By symmetry, double the sum on the positive half-line. The normal table gives
\begin{align*}
\P(1\le X<2)&=\Phi(2)-\Phi(1)\approx0.97725-0.84134=0.13591,\\
\P(0.1\le X<0.2)&\approx0.57926-0.53983=0.03943,\\
\P(0.01\le X<0.02)&\approx0.50798-0.50399=0.00399.
\end{align*}
The smaller intervals have total probability bounded above by
\[\P(0\le X<0.002)<\Phi(0.01)-\Phi(0)=0.00399.\]
(In fact $\Phi(0.002)-\Phi(0)\approx0.0008$, but the first bound is available from the table.) The upper tail is bounded by
\[\P(X\ge10)\le\P(X\ge3.99)\approx0.00003,\]
using the last table entry. The positive contribution is approximately $0.17933$ with error below $0.00402$. Thus the answer is about $0.35866$, with maximum error $0.00802$, giving about 36\% to accuracy 0.01. The source remarks that 37\% would also be accepted under its rough bound, and leaves the answer to accuracy $10^{-10}$ blank.
\subsection*{2. Exponential}
For $X\sim Exp(1)$, $F_X(t)=1-e^{-t}$ for $t\ge0$. The notes give
\[\P(X\le0.0002)\approx0.00020,\qquad\P(X>10)\approx0.00005.\]
They sum the central intervals and display
\[A\approx0.00010+0.00985+0.08610+0.23254=0.32859\approx33\%.\]
\noindent\textit{Transcription note: the source labels the first retained interval $[0.001,0.002)$ but prints its contribution as $0.00010$; the displayed approximation is retained as written.}
\subsection*{3. Ratio of uniforms}
Set $Y=U_1/U_2$ and $Z=U_2$. The inverse map is $U_1=YZ$, $U_2=Z$, with Jacobian $z$. Therefore
\[f_{Y,Z}(y,z)=z\one_{\{0\le yz\le1,\ 0\le z\le1\}}
=z\one_{\{y\ge0,\ 0\le z\le\min(1,1/y)\}}.\]
Hence
\[f_Y(y)=\int_0^{\min(1,1/y)}z\,dz=
\begin{cases}1/2,&0\le y\le1,\\1/(2y^2),&y>1,\\0,&y<0.\end{cases}\]
The source sketch is flat at height $1/2$ up to 1, then decreases as $1/(2y^2)$.
For the tails it gives $\P(Y<0.00002)=0.00001$ and, for $t\ge1$,
\[\P(Y>t)=\int_t^\infty\frac1{2y^2}\,dy=\frac1{2t},\qquad\P(Y>20000)\approx0.00003.\]
Summing the intervening intervals gives approximately $0.333325\approx33\%$. In this case the sum can also be evaluated exactly:
\begin{align*}
A&=\frac12\sum_{k=1}^\infty10^{-k}
+\sum_{k=0}^\infty\left(\frac1{2\cdot10^k}-\frac1{4\cdot10^k}\right)\\
&=\frac12\frac{1/10}{1-1/10}+\frac14\frac1{1-1/10}=\frac13.
\end{align*}
\subsection*{4. Powers of two}
The notes propose observing and proving that exactly one power of two with $k$ decimal digits begins with 1, for each $k$. Thus there are $k$ such powers up to $10^k$. The number of powers of two below $10^k$ is determined by $2^m<10^k$, or $m<k\log_2 10$. Their resulting bounds on the proportion are
\[\frac{k}{k\log_2 10+1}\le\frac{k}{\lceil k\log_2 10\rceil}\le\frac{k}{k\log_2 10},\]
and both outer terms tend to $1/\log_2 10=\log_{10}2$, the Benford probability.
\noindent\textit{Transcription note: the stated ``exactly one'' counting claim is reproduced from the handwritten argument; no replacement proof has been supplied.}
\section*{Problem 2: gamma ratio}
Set $Y=X_1/(X_1+X_2)$, $Z=X_1+X_2$. Then $X_1=YZ$, $X_2=Z-YZ$, and
\[J=\begin{vmatrix}z&y\\-z&1-y\end{vmatrix}=z.\]
The notes start from
\[f_{Y,Z}(y,z)=16(yz)^2e^{-2yz}(z-yz)^2e^{-2(z-yz)}z\one_{\{yz>0,z-yz>0\}},\]
then simplify it in the handwriting to
\[16y^2(1-y)^2z^3e^{-2z}\one_{\{z>0,0<y<1\}}.\]
\noindent\textit{Transcription note: this simplification has $z^3$ in the source, although the preceding product contains five powers of $z$. The subsequent workings below follow that displayed exponent.}

The notes write
\[f_Y(y)=16y^2(1-y)^2\int_0^\infty z^3e^{-2z}\,dz\]
and give three ways to evaluate the constant:
\begin{enumerate}
\item Integration by parts gives $\int_0^\infty z^3e^{-2z}\,dz=3/8$.
\item With $t=2z$, the integral is $\frac1{16}\int_0^\infty t^3e^{-t}dt=\frac1{16}\E X^3$ for $X\sim Exp(1)$. Its moment generating function is
\[M_X(\lambda)=\int_0^\infty e^{\lambda x}e^{-x}dx=\frac1{1-\lambda},\]
so $M_X'(0)=1$, $M_X''(0)=2$, $M_X'''(0)=6$, again giving $3/8$.
\item Since the expression must be a density on $(0,1)$, integrate $y^2(1-y)^2$ on $(0,1)$ to identify the normalising constant.
\end{enumerate}
For part 2, $Cor(X_1,X_2)=0$ because the variables are independent by assumption. The notes also conclude
\[Cor\left(\frac{X_1}{X_1+X_2},X_1+X_2\right)=0,\]
because the ratio and sum are independent: the displayed joint density factors into a function of $y$ and a function of $z$, including separate support indicators. This factorisation gives $f_{Y,Z}=f_Yf_Z$.
\section*{Problem 3: CLT approximation}
\begin{enumerate}
\item ``Standard CLT calculation.''
\item For $X_i\sim Ber(1/1000)$,
\[\mu=1/1000,\qquad\sigma=\sqrt{999/10^6}\approx10^{-1.5}.\]
The handwritten expression for the third moment is
\[\rho=\frac1{1000}\left(1-\frac1{1000}\right)^3+\frac{999}{1000}\left(0-\frac1{1000}\right)^3\approx10^{-3}.\]
It uses the Berry--Esseen error bound $\rho/(2\sigma^3\sqrt n)$. For $n=100$, the displayed estimate is
\[\frac\rho{2\sigma^3\sqrt{100}}\approx10^{-3}\cdot\frac12\cdot10^{4.5}\cdot\frac1{10}\approx1.58.\]
For desired errors 0.001, 0.01 and 0.1, the handwritten conclusions are respectively
\[\begin{array}{c|c|c}
\text{error}&\text{bound on }\sqrt n&\text{displayed bound on }n\\\hline
0.001&10^{1.5}\cdot10^3/2&10^9/2\\
0.01&10^2\cdot10^{1.5}/2&10^7/2\\
0.1&10\cdot10^{1.5}/2&10^5/2
\end{array}\]
\textit{The source writes the bounds with denominator 2 after squaring; those figures are transcribed unchanged.}

For comparison, when $p=1/2$, the notes use $\mu=1/2$, $\sigma=1/2$, $\rho=1/8$. The corresponding requirements are
\[\begin{array}{c|c|c}\text{error}&\sqrt n&n\\\hline0.001&>500&>250000\\0.01&>50&>2500\\0.1&>5&>25\end{array}\]
\end{enumerate}
\section*{Problem 4}
``Standard calculation.''
\end{document}
